%!TEX root = ../../thesis.tex

Utilising the measurement methods used previously, I search for a liquid that better replicates a biological impedance.
This work is of benefit to engineers of medical implant devices.
Having the ability to formulate a solution that mimics the electrical conditions inside a living being greatly simplifies testing.

It was discovered through collaboration with Saluda Medical, developers of stimulator implants, that a solution of one-tenth the standard solution of phosphate buffered saline is used as a test fluid.
This solution was the best substitute for a living spine that the engineers had.
It was held in large buckets in the laboratories for use whenever a quick test needed to be carried out.
Electrodes are submerged into these buckets in order to recreate the electrical conditions inside a human spine.

This was not the only way to simulate the impedance conditions inside a person.
As presented in~\cref{sect:sheep_measurements}, the use of anaesthetised sheep were also used.
A sheep's spine has roughly the same dimensions as a humans so it represents a sensible choice in terms of geometry.
However, the resources involved with conducting a live sheep trial were high; requiring use of a hospital operating theatre, surgeon veterinarian, and equipment.

Engineers at Saluda had no way of knowing how well those baths of saline represented a sheep's spine.
It was shown in \cref{sect:sheep_measurements} that the match between the two was weak.
With that knowledge, and the measurement techniques developed thus far, we now ask if it is possible to create a better match.
Such a match would reduce the number of surgical operations the test engineers might need to conduct, saving resources.


\section{Measurements}

This work is a very recent addition to the previous work at the time of writing.
The measurements were not conducted in a laboratory environment, so it is expected that some inaccuracies will exist.
These inaccuracies will be related to:
\begin{enumerate}
    \item a lack of temperature control, meaning temperatures fluctuated between 18 and 26 degrees;
    \item use of consumer grade scales, jeweller's scales for under 5 grams and kitchen above;
    \item use of consumer grade ingredients and distilled water.
\end{enumerate}
The measurement equipment used to conduct the measurements were the same as those used in the lab.

\subsection{Configuration}

The function generator, Agilent 33220A, and oscilloscope, Tektronix TPS 2024, are connected to the electrode array, St. Jude Medical - Octrode, as is shown in \cref{fig:creatingCSF_setup}.
A \SI{10}{\kilo\ohm} resistor was used to measure the current driven between electrodes eight and three.
It had a measured resistance of \SI{9.990}{\kilo\ohm}, as measured with a Fluke Digital Multimeter.

\begin{figure}
    \centering
    \includegraphics[scale=0.95]{content/pt2/graphics/CreatingCSF_setup}
    \caption{\label{fig:creatingCSF_setup}Diagram showing the measurement configuration used to measure the CPE response and resistivity of mixed solutions}
\end{figure}

\subsection{Procedure}

Each measurement run begin at the upper end of the frequency spectrum and proceeded towards the low frequency endpoint.
Starting with the higher frequencies offers a chance to confirm correct measurement set-up early in the measurement.
A sample at the lowest frequency can take over a minute to acquire, where the higher frequencies drop to under a second.
The same frequencies were chosen that were used to measure the situation in live sheep.
This allows easy comparison between the sheep data and the impedance of mixed solutions.

Measurements were completely automated via a Linux based computer running Python scripts.
These scripts controlled the output settings of the waveform generator and acquired the resulting waveforms from the oscilloscope.
The scripts had the ability to set the horizontal and vertical scales on the oscilloscope channels in order to ensure appropriate scales were used.
These scripts were written by myself and have been made freely available through my github account (\url{https://github.com/MarkHedleyJones/pyLabInstruments})

\begin{figure}
    \centering
    \includegraphics[width=\textwidth]{content/pt2/graphics/measurementFlowchart}
    \caption{\label{fig:creatingCSF_pythonFlowchart}Diagram showing the execution of the measurement script}
\end{figure}

The measurement procedure followed by the script is shown as a simplified flowchart in \cref{fig:creatingCSF_pythonFlowchart}.
The programme steps through each of the required frequencies, making sure the target voltage is developed across electrodes seven and eight, calculating the interface impedance.

The target voltage across the interface is \SI{20}{\milli\volt}.
This voltage was previously determined as a safe stimulus voltage in that it does not trigger Faradaic reactions at the electrode's surface.
Because the impedance of the interface changes with frequency, it is necessary to alter the output amplitude to keep the voltage across electrode seven and eight consistent.

This measurement configuration differs from the three channel method used in earlier measurements.
Those used channel three of the oscilloscope to watch the voltage developed between electrodes three and eight.

\subsection{Ingredients Tried}
To determine how certain additives affect the impedance of the interface, a wide range of ingredients were mixed and tested.
They were used in a semi-random order, chosen to maximise the usage of the quantity of distilled water available at the time.
\begin{itemize}
    \item Glycerol
    \item Methylated Spirits
    \item Sodium Carbonate
    \item Sodium Bicarbonate
    \item Table-salt (non-iodised)
    \item Gelatine
    \item Citric Acid
    \item Corn-flour
\end{itemize}

\subsection{Results}

\subsubsection*{Water, Corn-flour and Salt}

\begin{figure}
    \centering
    \includegraphics[width=\textwidth]{content/pt2/graphics/run12_600ml-distilledWater_ZVsF_graph_mag}
    \caption{\label{fig:recipe_water_mag}Plot of impedance magnitude versus frequency (log-log) for distilled water.}
\end{figure}

\begin{figure}
    \centering
    \includegraphics[width=\textwidth]{content/pt2/graphics/run12_600ml-distilledWater_ZVsF_graph_phase}
    \caption{\label{fig:recipe_water_phase}Plot of impedance phase versus frequency (log-log) for distilled water.}
\end{figure}

\begin{figure}
    \centering
    \includegraphics[width=\textwidth]{content/pt2/graphics/run14_170ml-distilledWater_250g-cornflour_ZVsF_graph_mag}
    \caption{\label{fig:recipe_cornflour_mag}Plot of impedance magnitude versus frequency (log-log) for \SI{250}{\gram} cornflour mixed with \SI{175}{\milli\litre} distilled water.}
\end{figure}

\begin{figure}
    \centering
    \includegraphics[width=\textwidth]{content/pt2/graphics/run14_170ml-distilledWater_250g-cornflour_ZVsF_graph_phase}
    \caption{\label{fig:recipe_cornflour_phase}Plot of impedance phase versus frequency (log-log) for \SI{250}{\gram} cornflour mixed with \SI{175}{\milli\litre} distilled water.}
\end{figure}

\begin{figure}
    \centering
    \includegraphics[width=\textwidth]{content/pt2/graphics/run14_175ml-distilledWater_250g-cornflour_1g9-salt_ZVsF_graph_mag}
    \caption{\label{fig:recipe_cornflour_salt_mag}Plot of impedance magnitude versus frequency (log-log) for \SI{250}{\gram} cornflour mixed with \SI{175}{\milli\litre} distilled water and \SI{1.9}{\gram} table-salt.}
\end{figure}

\begin{figure}
    \centering
    \includegraphics[width=\textwidth]{content/pt2/graphics/run14_175ml-distilledWater_250g-cornflour_1g9-salt_ZVsF_graph_phase}
    \caption{\label{fig:recipe_cornflour_salt_phase}Plot of impedance phase versus frequency (log-log) for \SI{250}{\gram} cornflour mixed with \SI{175}{\milli\litre} distilled water and \SI{1.9}{\gram} table-salt.}
\end{figure}

\begin{figure}
    \centering
    \includegraphics[width=\textwidth]{content/pt2/graphics/run14_180ml-distilledWater_250g-cornflour_1g9-salt_ZVsF_graph_mag}
    \caption{\label{fig:recipe_cornflour_salt_extraWater_mag}Plot of impedance magnitude versus frequency (log-log) for \SI{250}{\gram} cornflour mixed with \SI{180}{\milli\litre} distilled water and \SI{1.9}{\gram} table-salt.}
\end{figure}

\begin{figure}
    \centering
    \includegraphics[width=\textwidth]{content/pt2/graphics/run14_180ml-distilledWater_250g-cornflour_1g9-salt_ZVsF_graph_phase}
    \caption{\label{fig:recipe_cornflour_salt_extraWater_phase}Plot of impedance phase versus frequency (log-log) for \SI{250}{\gram} cornflour mixed with \SI{180}{\milli\litre} distilled water and \SI{1.9}{\gram} table-salt.}
\end{figure}

\Cref{fig:recipe_water_mag,fig:recipe_water_phase,fig:recipe_cornflour_phase,fig:recipe_cornflour_mag,fig:recipe_cornflour_phase,fig:recipe_cornflour_salt_mag,fig:recipe_cornflour_salt_phase,fig:recipe_cornflour_salt_extraWater_mag,fig:recipe_cornflour_salt_extraWater_phase} show a progression of measurements beginning with distilled water, then adding cornflour, then adding salt, and finally adding a very small amount of extra water.

\subsubsection{Improved Water, Corn-flour, and Salt}

\begin{figure}
    \centering
    \includegraphics[width=\textwidth]{content/pt2/graphics/run12_190ml-distilledWater_190g-cornflour_0g858-salt_ZVsF_graph_mag}
    \caption{\label{fig:recipe_cornflour_salt_extraWater_mag_improved}Plot of impedance magnitude versus frequency (log-log) for \SI{190}{\gram} cornflour mixed with \SI{190}{\milli\litre} distilled water and \SI{0.858}{\gram} table-salt.}
\end{figure}

\begin{figure}
    \centering
    \includegraphics[width=\textwidth]{content/pt2/graphics/run12_190ml-distilledWater_190g-cornflour_0g858-salt_ZVsF_graph_phase}
    \caption{\label{fig:recipe_cornflour_salt_extraWater_phase_improved}Plot of impedance phase versus frequency (log-log) for \SI{190}{\gram} cornflour mixed with \SI{190}{\milli\litre} distilled water and \SI{0.858}{\gram} table-salt.}
\end{figure}

\Cref{fig:recipe_cornflour_salt_extraWater_mag_improved,fig:recipe_cornflour_salt_extraWater_phase_improved} show the results of mixing the same set of ingredients as the previous measurement set, but with the closest fit to sheep spine obtained.


\subsection{Discussion}

These results show that it is possible to alter the CPE response and the bulk conductivity of the solution independently of one another.

Each of the measurements that contained salt had a tendency to ``fall-over'' at low frequency.
The salt
